\documentclass[11pt]{article}
\usepackage[margin=1in]{geometry}
\usepackage{amsmath,amssymb,booktabs}
\usepackage{hyperref}
\title{Bayesian Linear Models as Priors for Safe Reinforcement Learning in Robotics}
\author{Syed Hafiz Ali}
\date{Research manuscript in preparation}
\begin{document}
\maketitle
\begin{abstract}
This manuscript specifies a research direction for using uncertainty-aware Bayesian estimates to inform intervention decisions in autonomous control. The companion implementation currently demonstrates a scalar Normal--Normal posterior update and interval-based ALLOW, VERIFY, or BLOCK gate. It is not a complete Bayesian reinforcement-learning algorithm and provides no formal safety guarantee.
\end{abstract}
\section{Motivation}
Safe reinforcement learning treats constraints on learning and deployment as part of the problem, not merely as a reward penalty \cite{garcia2015safe,berkenkamp2017safe}.
An autonomous controller should not treat a point prediction as certainty. When predicted risk is close to a decision threshold, uncertainty can justify evidence gathering or human review rather than immediate action.
\section{Scalar model}
For a scalar latent hazard parameter with Gaussian prior $\theta\sim\mathcal{N}(\mu_0,\sigma_0^2)$ and observations with known Gaussian noise variance, the posterior remains Gaussian. The prototype implements the conjugate update and forms an interval using a configurable multiplier.
\section{Decision rule}
For an interval $[\ell,u]$ and threshold $\tau$, authorize only when $u\leq\tau$, block when $\ell>\tau$, and otherwise request verification. This is a decision heuristic; interval calibration and operational costs must be evaluated.
\section{Research questions}
How does uncertainty-aware gating compare with point-estimate thresholding? How sensitive is it to prior misspecification, observation noise, and distribution shift? What is the cost of verification relative to false authorization?
\section{Evaluation design}
Use labelled synthetic hazard distributions and matched test cases. Compare point estimates, fixed-margin thresholds, and posterior intervals. Report false-allow, false-block, verification, coverage, calibration, and intervention rates across uncertainty regimes.
\section{Results}
No empirical safety or calibration results are claimed. Results must be generated from preregistered configurations and reported with uncertainty intervals.
\section{Limitations}
The current model is scalar, assumes known observation variance, and does not represent sequential state transitions, policy learning, correlated observations, or adversarial dynamics. A credible interval is not a safety certificate.
\section{Conclusion}
Bayesian uncertainty can inform when to defer, but safe control requires calibrated models, robust evaluation, and system-level analysis beyond the present prototype.
\begin{thebibliography}{9}
\bibitem{garcia2015safe} J. García and F. Fernández. A Comprehensive Survey on Safe Reinforcement Learning. \textit{Journal of Machine Learning Research}, 16(42):1437--1480, 2015. \url{https://www.jmlr.org/papers/v16/garcia15a.html}.
\bibitem{berkenkamp2017safe} F. Berkenkamp, M. Turchetta, A. P. Schoellig, and A. Krause. Safe Model-based Reinforcement Learning with Stability Guarantees. \textit{Advances in Neural Information Processing Systems}, 30, 2017. \url{https://arxiv.org/abs/1705.08551}.
\bibitem{recht2019tour} B. Recht. A Tour of Reinforcement Learning: The View from Continuous Control. \textit{Annual Review of Control, Robotics, and Autonomous Systems}, 2:253--279, 2019. DOI: \href{https://doi.org/10.1146/annurev-control-053018-023825}{10.1146/annurev-control-053018-023825}.
\end{thebibliography}
\end{document}
